\section{Analisi e Contromisure del Worm}\label{sec:Contromisure}
Il riferimento principale per la risposta all'incidente è l’advisory CA-2000-04, \emph{Love Letter Worm}, pubblicato dal CERT Coordination Center (CERT/CC) il 4 maggio 2000, che descriveva il worm e forniva raccomandazioni specifiche \cite{cert2000}. Le misure si possono organizzare in quattro categorie (prevenzione, rilevamento, contenimento e ripristino), che riprendono le fasi della gestione degli incidenti descritte dal NIST \cite{cichonski2012}.

\subsection{Prevenzione}
Il CERT raccomandava di impedire l’esecuzione degli script disabilitando Windows Script Host (WSH) e, ove opportuno, lo scripting attivo in Internet Explorer, oltre a mantenere aggiornati gli antivirus per riconoscere le varianti note \cite{cert2000}.

Sul piano umano, invitava gli utenti a prestare attenzione (\emph{``exercise caution''}): verificare la provenienza dei messaggi e non aprire allegati inattesi \cite{cert2000, bishop2000}.

\subsection{Rilevamento}
Nel 2000 il rilevamento si basava su antivirus aggiornati, scansioni dei file e indicatori riconoscibili, come l’allegato \texttt{LOVE-LETTER-FOR-YOU.TXT.vbs} e l'oggetto \emph{ILOVEYOU} \cite{cert2000, bishop2000}. La sua efficacia dipendeva dalla rapidità con cui gli avvisi del CERT e dei produttori di antivirus arrivavano agli amministratori e venivano applicati \cite{cert2000, gao2000iloveyou}.

\subsection{Contenimento}
Il contenimento richiedeva rapidità, poiché il worm si propagava tramite Outlook (Sezione~\ref{subsec:fase2}), IRC e cartelle condivise \cite{cert2000, bishop2000}. Le azioni principali erano: il blocco dei messaggi infetti, ovvero filtrare sui server di posta i messaggi con oggetto ``ILOVEYOU'' e gli allegati \texttt{.vbs}, l'interruzione dei canali di diffusione alternativi, quindi disabilitare la ricezione automatica dei file via DCC nei client IRC \cite{cert2000}, l'isolamento delle macchine compromesse, scollegando dalla rete i sistemi infetti per impedire ulteriori invii \cite{cichonski2012} e, infine, il coordinamento tra amministratori e organizazioni, ossia la condivisione tempestiva di indicatori e misure di blocco tra organizzazioni.

\subsection{Ripristino}
Il ripristino prevedeva l’uso degli strumenti di pulizia dei produttori di antivirus, il recupero dei file da copie integre e la verifica che non restassero componenti o configurazioni capaci di riattivare l'infezione, prima di reinserire il sistema in rete \cite{cichonski2012, nelson2025}. L’incidente mostrò così l’importanza di copie di sicurezza affidabili.

\subsection{Ulteriori considerazioni tecniche}
Matt Bishop, nella sua analisi del worm \cite{bishop2000}, evidenzia i limiti dei filtri basati su elementi del messaggio come l’oggetto, che può essere facilmente modificato. Propone invece approcci strutturali: il \emph{sandboxing}, cioè l'esecuzione dei programmi in un ambiente isolato e limitato, e il principio del minimo privilegio, per cui ogni programma dispone solo degli accessi strettamente necessari. Così anche un allegato malevolo eseguito avrebbe poche possibilità di modificare file o usare altri programmi. Suggerisce inoltre di individuare comportamenti anomali, come l’invio massivo di email o modifiche inattese al registro, anziché affidarsi solo alle firme antivirus.

\subsection{Oggi}
Nel 2000 la risposta si basava su avvisi, aggiornamento delle firme e applicazione manuale delle contromisure \cite{cert2000, gao2000iloveyou}. Oggi le organizzazioni dispongono di strumenti più integrati, in un approccio basato sulla gestione continua del rischio \cite{nelson2025, nist2024csf}:

\begin{itemize}
    \item \textbf{Prevenzione:} i gateway email bloccano gli allegati pericolosi e li analizzano in sandbox, mentre l'\emph{application control} applica il minimo privilegio proposto da Bishop \cite{bishop2000, nist2024csf}.
    \item \textbf{Rilevamento:} gli strumenti di monitoraggio degli endpoint e i SIEM (\emph{Security Information and Event Management}) analizzano e correlano le attività dei dispositivi, individuando comportamenti anomali anche senza una firma nota \cite{nelson2025, nist2024csf}.
    \item \textbf{Contenimento:} le piattaforme di sicurezza isolano automaticamente gli endpoint sospetti e bloccano i messaggi dannosi, mentre la segmentazione della rete limita l'estensione di un'infezione \cite{nelson2025, nist2024csf}.
    \item \textbf{Ripristino:} backup verificati e immutabili e procedure documentate riducono i tempi di inattività e la perdita di dati \cite{nelson2025, nist2024csf}.
    \item \textbf{Formazione:} sensibilizzazione, simulazioni di \emph{phishing} ed esercitazioni preparano utenti e tecnici a riconoscere e gestire una minaccia \cite{nelson2025, nist2024csf}.
\end{itemize}

L'approccio diventa così più proattivo, ma il rischio non è eliminato: configurazioni errate, strumenti non aggiornati e tecniche di \emph{social engineering} sempre più sofisticate possono ancora far riuscire un attacco \cite{nelson2025, nist2024csf}.
